% =============================================================================
\section{Conclusion}\label{sec:conclusion}
% =============================================================================

We presented a forecast-to-decision pipeline for NHS critical-care
surge planning that couples a per-region physics-informed forecaster to
a tail-risk robust allocation linear program. With a leakage-free
per-origin recalibration, PinnGRU is competitive across the horizons
that drive surge decisions, cutting $14$-day RMSE by $25\%$ over
re-fitted ARIMA and $42\%$ over a neural GRU. On the decision side the
gain is structural: once the budget binds, every budget-exhausting
policy ties on coverage and the robust formulation wins by routing the
same surge with the least patient transfer. Principal caveats are
detour-scaled travel times, well-mixed populations, scaled cost units,
and a single realised origin. A deeper limitation couples the modules:
the per-region forecaster and co-monotone scenarios over-state the joint
tail and understate transfers, which pay off most under
\emph{asynchronous} regional peaks. Correlated demand scenarios from the
joint forecast distribution, and a trust-level instance where exact
solution grows expensive enough to warrant metaheuristics, are the
natural next steps.
